\documentclass[pdflatex,sn-mathphys-num]{sn-jnl}
\usepackage{graphicx}
\usepackage{multirow}
\usepackage{amsmath,amssymb,amsfonts}
\usepackage{amsthm}
\usepackage{mathrsfs}
\usepackage[title]{appendix}
\usepackage{xcolor}
\usepackage{textcomp}
\usepackage{manyfoot}
\usepackage{booktabs}
\usepackage{listings}
\usepackage{longtable}
\usepackage{calc}
\usepackage{array}
\usepackage{tabularx}
\usepackage{bookmark}
\lstset{basicstyle=\footnotesize\ttfamily,breaklines=true,breakatwhitespace=true,columns=fullflexible,keepspaces=true,showstringspaces=false,xleftmargin=0pt,breakindent=1.5em,postbreak=\mbox{\textcolor{gray}{$\hookrightarrow$}\space},frame=none,aboveskip=6pt,belowskip=4pt}
\usepackage{newunicodechar}
\newunicodechar{×}{\ensuremath{\times}}
\newunicodechar{→}{\ensuremath{\rightarrow}}
\newunicodechar{←}{\ensuremath{\leftarrow}}
\newunicodechar{≈}{\ensuremath{\approx}}
\newunicodechar{≥}{\ensuremath{\geq}}
\newunicodechar{≤}{\ensuremath{\leq}}
\newunicodechar{≠}{\ensuremath{\neq}}
\newunicodechar{ρ}{\ensuremath{\rho}}
\newunicodechar{Φ}{\ensuremath{\Phi}}
\newunicodechar{φ}{\ensuremath{\phi}}
\newunicodechar{α}{\ensuremath{\alpha}}
\newunicodechar{β}{\ensuremath{\beta}}
\newunicodechar{χ}{\ensuremath{\chi}}
\newunicodechar{θ}{\ensuremath{\theta}}
\newunicodechar{σ}{\ensuremath{\sigma}}
\newunicodechar{μ}{\ensuremath{\mu}}
\newunicodechar{λ}{\ensuremath{\lambda}}
\newunicodechar{π}{\ensuremath{\pi}}
\newunicodechar{ε}{\ensuremath{\epsilon}}
\newunicodechar{δ}{\ensuremath{\delta}}
\newunicodechar{γ}{\ensuremath{\gamma}}
\newunicodechar{η}{\ensuremath{\eta}}
\newunicodechar{Σ}{\ensuremath{\Sigma}}
\newunicodechar{Δ}{\ensuremath{\Delta}}
\newunicodechar{±}{\ensuremath{\pm}}
\newunicodechar{∞}{\ensuremath{\infty}}
\newunicodechar{∈}{\ensuremath{\in}}
\newunicodechar{·}{\ensuremath{\cdot}}
\newunicodechar{°}{\textdegree}
\newunicodechar{∑}{\ensuremath{\sum}}
\newunicodechar{²}{\textsuperscript{2}}
\newunicodechar{³}{\textsuperscript{3}}
\newunicodechar{¹}{\textsuperscript{1}}
\newunicodechar{–}{--}
\newunicodechar{—}{---}
\newunicodechar{ł}{\l}
\newunicodechar{Ł}{\L}
\newunicodechar{ż}{\.z}
\newunicodechar{Ż}{\.Z}
\newunicodechar{ó}{\'o}
\newunicodechar{ś}{\'s}
\newunicodechar{ą}{\k a}
\newunicodechar{ę}{\k e}
\newunicodechar{ć}{\'c}
\newunicodechar{ń}{\'n}
\newunicodechar{ź}{\'z}
\newunicodechar{’}{'}
\newunicodechar{‘}{`}
\newunicodechar{“}{``}
\newunicodechar{”}{''}
\newunicodechar{…}{\ldots}
\providecommand{\tightlist}{\setlength{\itemsep}{0pt}\setlength{\parskip}{0pt}}
\providecommand{\passthrough}[1]{#1}
\providecommand{\pandocbounded}[1]{#1}
\hypersetup{hypertexnames=false}
\raggedbottom
\begin{document}
\title[Variance Components of LLM Non-Determinism]{Where Does the Noise Come From? A Variance-Components Decomposition of Non-Determinism in LLM Brand Answers}

\author*[1,2]{\fnm{Dmitrij} \sur{\.Zatuchin}}\email{dmitrij.zatuchin@eek.ee}

\affil*[1]{\orgdiv{Department of Information Technologies}, \orgname{Estonian Entrepreneurship University of Applied Sciences (EUAS)}, \orgaddress{\city{Tallinn}, \country{Estonia}}}
\affil[2]{\orgname{Rankfor.AI}, \orgaddress{\city{Wroc\l{}aw}, \country{Poland}; \city{Tallinn}, \country{Estonia}}}

\abstract{Teams that measure whether large language models (LLMs) recommend a brand face a reproducibility problem: ask the same question twice and the answer moves. Practice has converged on resampling each prompt a small fixed number of times, commonly five, then averaging. That convention treats within-prompt resampling as the source of the noise. A measured brand score can move for at least four separable reasons: within-prompt resampling, prompt paraphrase, model identity, and query language. These cost different amounts to sample and enter the precision of a brand estimate through different denominators. We specify a crossed random-effects (generalizability-theory) decomposition that partitions the total variance of a response-level brand outcome into those four sources, and we embed the components in a decision-study allocation that returns how many repeats, paraphrases, models, and languages to buy for a target reliability. We apply the method to a fully crossed corpus of 12,933 LLM responses about 20 Central and Eastern European brands in 8 languages across 3 models (GPT-5.2 and Gemini 3 Flash queried in parametric mode, Perplexity in grounded retrieval mode), with 15 prompts per brand-language-model cell and a stability subset of 1,435 cells resampled about five times each (7,173 responses). The modeled outcome is per-response multilingual sentiment polarity toward the brand, of which 91.9\% of responses score exactly neutral; the outcome is a continuous response-level score, not a recommendation or ranking indicator, and the recommendation-indicator refit is named future work (Section 5.4). We report the fitted variance components, their intraclass correlations, the precision-versus-budget frontier, and the allocation rule from converged REML fits. Query language is the largest systematic facet: it accounts for 26.5\% of the variance of one response, against 1.5\% for brand identity (ICC 0.0146), so a single AI answer carries almost no brand-discriminating signal. On the resampled stability subset, once a cell term isolates pure within-prompt resampling, the resampling component is 34.8\% of variance and the brand-in-context interaction previously hidden inside the residual is 29.6\%. In the object-by-facet fit the brand's context-free true score is 0.7\%, and the only sizeable brand-by-facet term is brand-by-language at 8.6\%, a measurable bilingual penalty; brand-by-model and brand-by-prompt are near zero. The decision study confirms the structural prediction: per unit of query budget, adding languages and models reduces relative-error variance far more than adding repeats, and a repeat past the fifth reduces it by only 0.0003. Brand-ranking reliability stays low across the whole design grid, near 0.01 for a single answer and about 0.36 even at the full crossed 8-language, 3-model, 15-paraphrase design, so reliability is bought by spreading across languages and models, not by repeating one prompt. Cluster-bootstrap confidence intervals on the components are the remaining v2 finalization.}

\keywords{generalizability theory; variance components; intraclass correlation; large language models; measurement reproducibility; sampling design}

\maketitle

\section{Introduction}\label{introduction}

An LLM answering ``which parcel carrier should I use in Poland?'' does not return a fixed answer. Ask it again and the brands, their order, and the tone can change. Anyone who measures this, whether for academic study of AI-mediated brand reputation or for commercial measurement, confronts the same question before any finding is trustworthy: how many times must you ask, and in how many ways, before the measured answer stops being an artifact of the sampling procedure.

The common answer is a convention. Query each prompt a fixed small number of times, often five, average the outcome, and report the mean. Our own earlier work adopted five iterations as an exploratory tier \cite{ref19}. Independent industry measurement reports the raw instability that motivates the convention: one vendor found only 2.2\% of cited sources identical across three repeated runs, with week-to-week citation shifts of 56 to 74\% (Indig, via Search Engine Land, June 2026 \cite{ref9}); another estimated that a stable position estimate needs about 10 runs for a quick read and anywhere from 7 to 135 runs for a tight confidence interval (Graphite, May 2026, 200 prompts by 400 responses \cite{ref8}). These figures treat the number of repeated runs as the lever, and they place all of the instability in resampling.

That framing has a blind spot. A measured brand score can move for at least four distinct reasons. It moves because the decoder samples a different continuation on an identical prompt (\textbf{within-prompt resampling}). It moves because the same question, phrased differently, pulls a different answer (\textbf{prompt paraphrase}). It moves because a different model was asked (\textbf{model identity}). It moves because the question was asked in a different language (\textbf{query language}). These four cost different amounts to sample, and they enter the precision of a brand estimate through different denominators. A study that spends its entire query budget on repeats may be paying down one source of error while leaving the others untouched. Whether that is what happens, and by how much, is an empirical question about the sizes of the four sources, and this paper answers it for one corpus and one outcome.

This paper does three things.

\begin{enumerate}
\def\labelenumi{\arabic{enumi}.}
\item
  \textbf{Decomposition.} We specify a crossed random-effects model that partitions the total variance of a response-level brand outcome into the four sources above plus their higher-order interactions, each reported as an intraclass correlation (ICC).
\item
  \textbf{Allocation.} We embed the components in the generalizability-theory decision-study machinery and derive a precision-versus-query-budget frontier, together with a rule that returns the number of repeats, paraphrases, models, and languages required to reach a target reliability at least cost.
\item
  \textbf{The empirical question.} We ask whether repeats beyond the common five add information once paraphrase and model variance are in the design. The structure of the error variance predicts that they add less per query than the other facets, and the fitted components confirm it: a repeat past the fifth reduces relative-error variance by about 0.0003, while spreading the same budget across languages or models reduces it several times more.
\end{enumerate}

The contribution is the method, the estimand specification, the structural argument that does not depend on the fitted numbers, and the fitted decomposition itself. The variance components, the precision-versus-budget frontier, and the allocation are computed on the corpus described in Section 3 and reported in Section 5; cluster-bootstrap confidence intervals on the components are the one remaining finalization step. Section 2 places the decomposition against two neighboring lines of our own work. Section 3 describes the corpora. Section 4 states the models and the allocation rule. Sections 5 to 8 give results, discussion, limitations, and the ethics and data-availability statements.

\section{Related work}\label{related-work}

\textbf{Generalizability theory.} The statistical apparatus is standard. Generalizability theory \cite{ref4,ref6,ref16} was built to answer the allocation question, how many raters, items, and occasions a measurement needs, by decomposing an observed score into variance components for an object of measurement and a set of crossed or nested facets, then using those components in a decision study to project the reliability of any proposed design. Variance-component estimation by restricted maximum likelihood (REML) in crossed, unbalanced designs is well established \cite{ref14}, and its modern software form is the linear mixed model \cite{ref3,ref15}. Variance partitioning into ICCs is likewise standard \cite{ref10}. The contribution here is the identification of the four LLM-specific facets, the placement of resampling as the facet that a fixed budget already suppresses through averaging, and the application of the decision-study projection to LLM query budgets.

\textbf{Measuring LLM brand visibility.} A body of applied measurement, much of it from commercial vendors, reports raw run-to-run instability in LLM answers and citations \cite{ref8,ref9,ref12,ref13}. These reports establish that the noise is large and that sampling matters. They stop short of separating the noise into sources or of turning the separation into a spending rule. Peer-reviewed work anchors the resampling facet: Ouyang, Zhang, Harman and Wang \cite{ref11} show that ChatGPT's outputs vary substantially across repeated identical code-generation requests, and Atil et al.~\cite{ref1} find that non-determinism persists across runs of hosted LLM APIs configured to be deterministic, with accuracy swings across ten runs that survive even at temperature 0. The floating-point and batch-size mechanism behind this residual non-determinism is documented in ``Understanding and Mitigating Numerical Sources of Nondeterminism in LLM Inference'' \cite{ref17}. We cite the vendor measurements as the current public evidence of the run-to-run instability these mechanisms produce in the wild, each with its sample and date.

\textbf{Two neighboring studies, and what is new here.} This work sits next to two prior papers and is distinct from both.

\begin{itemize}
\item
  The \emph{category-ownership map} \cite{ref18} (3,750 responses, 50 brands, 5 industries, 250 category queries, 3 models; dataset DOI 10.5281/zenodo.20788142) measures how recommendation share is distributed across brands within a category, using the Gini coefficient and entropy (reported mean Gini 0.28, 95\% CI 0.16 to 0.41). Those are concentration indices. They describe whether a category is a monopoly or a scramble. They do not describe why a given brand's measured share wobbles between runs, or which sampling facet is responsible. Concentration is a property of the answer; the present decomposition is a property of the measurement error around the answer.
\item
  The \emph{dice-roll meta-methodology} \cite{ref19} (Monte Carlo power simulation and generalizability analysis over roughly 190,000 observations) sets power tiers on a single facet, the number of iterations: exploratory n=5 (G=0.58), confirmatory n=10 (G=0.74), rigorous n=15 (G=0.81). That work asks how many repeats buy adequate power. The present work asks the crossed question: given four facets that all inject variance, how should a fixed query budget be split among repeats, paraphrases, models, and languages. The two are complementary. The dice-roll tiers are the special case of the allocation rule below when languages, models, and paraphrases are held at one.
\end{itemize}

\section{Data}\label{data}

We use three corpora. The first is the load-bearing, fully crossed corpus for the decomposition in Section 5. The other two are secondary corpora, described here because the same decomposition applies to them and because they extend the design to a corporate-taxonomy and a single-brand setting. Their decompositions are named as generalization fits and are not run in this version (Section 5.4).

\subsection{Primary corpus: the CEE cross-language corpus}\label{primary-corpus-the-cee-cross-language-corpus}

The primary corpus is the Central and Eastern European (CEE) arm of a cross-language AI brand-reputation dataset \cite{ref20} (dataset DOI 10.5281/zenodo.20794390). Twenty brands headquartered in CEE markets were probed in eight languages (English, Polish, Czech, Slovak, Hungarian, Lithuanian, Latvian, Estonian) across three models, using fifteen prompts organized into five categories of three prompts each. The eight languages span three CEE language families, Slavic (Polish, Czech, Slovak), Uralic (Hungarian, Estonian), and Baltic (Lithuanian, Latvian), with English as a Germanic reference; the corpus therefore carries four language-family labels. The three models are GPT-5.2 and Gemini 3 Flash (\texttt{gemini-3-flash-preview}), both queried in parametric mode without retrieval, and Perplexity, queried in grounded retrieval mode. All three were run at temperature 0.3. The brands are Allegro, Alza.cz, Bolt, CCC, CD Projekt, DPD, ESET, InPost, Kiwi.com, Kofola, MOL Group, Notino, Packeta, Reserved, Revolut, Vinted, Wise, Wizz Air, Wolt and mBank, tiered by recognition (pan-CEE, regional, domestic). Collection ran in a single window in spring 2026.

The design makes the decomposition possible. Brand, language, model and prompt are crossed, so the variance attributable to each can be separated from the others. The three prompts in Category D (stability probes, ``Describe {[}Brand{]} in detail'' and two paraphrases) were repeated five times per brand-language-model cell; every other cell was queried once. Category D is therefore the only place in the corpus where pure within-prompt resampling variance is identified. The fully crossed design specifies 12,960 target responses: 5,760 single-shot cells (categories A, B, C, E) plus 7,200 across the 1,440 Category D cells at five iterations each. Collection issued 14,400 API calls; the surplus over the design is retries and calls that returned errors or empty responses. Of the 12,960 target responses, 12,933 are usable (99.8\% of the design; the 27-response shortfall comes from five Category D cells that produced no usable response and two that returned four iterations instead of five). Table 1 gives the counts, all verified against the response-level table.

\begin{table}[htbp]
\caption{Primary corpus structure (n = 12,933 responses).}
\small
\begin{tabularx}{\textwidth}{@{}>{\raggedright\arraybackslash}X>{\raggedright\arraybackslash}X>{\raggedright\arraybackslash}X@{}}
\toprule
Facet & Levels & Note \\
\midrule

Brand (object of measurement) & 20 & CEE-headquartered, 5 industries, 3 recognition tiers \\
Query language & 8 & Slavic (PL, CZ, SK), Uralic (HU, ET), Baltic (LT, LV), English (reference) \\
Model & 3 & GPT-5.2 (parametric), Gemini 3 Flash (parametric), Perplexity (grounded RAG); temperature 0.3 \\
Prompt & 15 & 5 categories (A open reputation, B source elicitation, C competitive, D stability, E industry) of 3 each \\
Resampling (iteration) & 1, or \textasciitilde5 on the Category D subset & 1,435 cells (1,433 at 5 iterations, 2 at 4) = 7,173 responses \\
Single-shot responses (categories A, B, C, E) & 20 x 8 x 3 x 12 & 5,760 responses \\
\textbf{Total usable responses} & \textbf{12,933} & 7,173 stability + 5,760 single-shot \\
\bottomrule
\end{tabularx}
\end{table}

\textbf{Outcome.} The response-level outcome is the sentiment polarity of each response toward the target brand, in {[}-1, 1{]}, scored by a multilingual transformer, \texttt{cardiffnlp/twitter-xlm-roberta-base-sentiment-multilingual}, an XLM-RoBERTa sentiment classifier (Conneau et al.~\cite{ref5} for the base encoder; Barbieri, Anke and Camacho-Collados \cite{ref2} for the multilingual sentiment fine-tune). Sentiment is the continuous per-response signal available across every language and model in the corpus. Two properties attach to it and are carried through the analysis. First, the outcome is strongly zero-inflated: the median is 0, the mean is -0.0205, and 91.9\% of responses score exactly neutral, so a Gaussian mixed model treats a near-degenerate outcome and the components should be read as a first-order decomposition. Second, the scorer is noisy in absolute terms on long, list-structured answers. The decomposition therefore concerns relative variation across facets, and not absolute calibration. A recommendation-indicator outcome (whether the brand was named, and in what rank) is the natural second outcome and is the pre-registered extension named in Section 5.4.

For context on the outcome's structure, prior analysis of these same data \cite{ref20} reports cross-language sentiment divergence between English and each local language: 16 of the 20 brands diverge by more than 0.15 from their English sentiment in all seven local languages, 18 of 20 diverge in at least one, and only MOL Group and mBank show no divergence above 0.15 in any local language. The largest single divergences reach roughly 0.33 (for example Bolt in Polish). This is a divergence measure on aggregated per-brand-per-language sentiment, not a variance component; it establishes that the language facet carries a large signal on this outcome, and it is one input to the interpretation of the fitted decomposition.

\subsection{Secondary corpus: corporate taxonomy (six large corporates)}\label{secondary-corpus-corporate-taxonomy-six-large-corporates}

The second corpus applies the dice-roll protocol to six large Polish and European corporates (PGE, Polpharma, Orlen, Holcim, TotalEnergies, V-ZUG) across 195 iterations \cite{ref21}. It documents three qualitative visibility patterns (invisible, brand-dependent, topic-locked) and is included here as a corporate-taxonomy setting to which the same crossed decomposition applies. Its variance-components fit is a generalization fit and is not run in this version (Section 5.4).

\subsection{Secondary corpus: V-ZUG single-brand deep probe}\label{secondary-corpus-v-zug-single-brand-deep-probe}

The third corpus is a single-brand unbranded-query probe of the Swiss premium appliance maker V-ZUG, 375 queries (25 prompts x 5 iterations x 3 models), reporting 0\% visibility on one model (Grok) and 3.2\% on another (Gemini), with the incumbent Miele appearing 5 to 26 times more often \cite{ref22}. It is the single-brand, high-replication end of the design space, with five repeats, one language, three models, and twenty-five prompts, and is included to show where a repeat-heavy, paraphrase-light design would sit on the frontier of Section 5.3. Its variance-components fit is a generalization fit and is not run in this version (Section 5.4).

\section{Method}\label{method}

\subsection{The crossed variance-components model}\label{the-crossed-variance-components-model}

Let \(y_{b\ell m p i}\) be the outcome for brand \(b\), language \(\ell\), model \(m\), prompt \(p\), iteration \(i\). Treating brand as the object of measurement and language, model and prompt as crossed random facets, the response is

\[
y_{b\ell m p i} = \mu + \beta_b + \lambda_\ell + \gamma_m + \pi_p + (\text{interactions}) + \varepsilon_{b\ell m p i},
\]

with every term after \(\mu\) an independent, mean-zero random effect: \(\beta_b \sim N(0,\sigma^2_{\text{brand}})\), \(\lambda_\ell \sim N(0,\sigma^2_{\text{lang}})\), \(\gamma_m \sim N(0,\sigma^2_{\text{model}})\), \(\pi_p \sim N(0,\sigma^2_{\text{prompt}})\), the interaction terms as noted below, and the residual \(\varepsilon \sim N(0,\sigma^2_{\text{resid}})\). The residual is pure within-prompt resampling only where iterations replicate the same cell, which is why the resampling component is identified on the Category D subset and not on the single-shot cells.

We fit three nested specifications by REML.

\begin{itemize}
\item
  \textbf{M1, main-effects decomposition (full corpus, n = 12,933).} Random intercepts for brand, language, model, prompt; residual. On the single-shot cells the residual absorbs all interactions together with resampling, so M1 answers where the variance of one response comes from but does not isolate resampling.
\item
  \textbf{M2, resampling-isolating decomposition (Category D subset, n = 7,173).} Random intercepts for brand, language, model, prompt, plus a random intercept for the full brand x language x model x prompt cell. With the cell term present, every between-cell interaction is captured by the cell effect and the residual becomes pure within-prompt resampling. We validate this by comparing the fitted residual to the model-free mean within-cell variance computed directly from the replicated cells.
\item
  \textbf{M3, object-by-facet decomposition (Category D subset, n = 7,173).} M2 plus the object-by-facet interactions brand x language, brand x model, brand x prompt, which are the error terms that matter for ranking brands. The facet main effects matter for scoring a brand on a fixed scale. M3 supplies the components for the decision-study allocation.
\end{itemize}

The ICC for a source is its variance divided by the total. Confidence intervals come from a nonparametric cluster bootstrap that resamples the 20 brands with replacement, refits each specification, and takes the 2.5th and 97.5th percentiles over converged replicates; resampling the object of measurement is the honest bootstrap unit for generalizing to new brands, and with 20 brands the intervals are wide by construction \cite{ref7}. The intervals are the one remaining finalization step (Section 5.4); this version reports the point estimates.

\subsection{The decision study and the allocation rule}\label{the-decision-study-and-the-allocation-rule}

For a design that averages a brand's outcome over \(n_L\) languages, \(n_M\) models, \(n_P\) paraphrases and \(n_R\) repeats, generalizability theory gives the error variance by dividing each component by the product of the sample sizes of the facets it indexes. Writing the M3 components as \(\sigma^2_{p}\) (brand), \(\sigma^2_{pL},\sigma^2_{pM},\sigma^2_{pP}\) (brand-by-facet), \(\sigma^2_{\text{cell}}\) (highest-order brand-inclusive interaction) and \(\sigma^2_{e}\) (resampling), the relative error variance for ranking brands is

\[
\sigma^2_\delta = \frac{\sigma^2_{pL}}{n_L} + \frac{\sigma^2_{pM}}{n_M} + \frac{\sigma^2_{pP}}{n_P} + \frac{\sigma^2_{\text{cell}}}{n_L n_M n_P} + \frac{\sigma^2_{e}}{n_L n_M n_P n_R},
\]

and the absolute error variance for scoring a brand on a fixed scale adds the facet main effects \(\sigma^2_L/n_L + \sigma^2_M/n_M + \sigma^2_P/n_P\). The generalizability coefficient is \(E\rho^2 = \sigma^2_p / (\sigma^2_p + \sigma^2_\delta)\) and its absolute counterpart is \(\Phi = \sigma^2_p / (\sigma^2_p + \sigma^2_\Delta)\).

The structure of \(\sigma^2_\delta\) carries the standing argument, independent of the fitted values. The resampling term is divided by \(n_L n_M n_P n_R\), the full query count. The language term is divided by only \(n_L\). So the marginal reduction from one more repeat is bounded by \(\sigma^2_e\) over an already-large product, while the marginal reduction from one more language is \(\sigma^2_{pL}\) over a small one. For any positive component values, the repeat facet has the fastest-decaying marginal value. Whether the resulting gap between facets is small or large for this outcome depends on the fitted magnitudes, which Section 5 reports. The allocation rule is a discrete optimization: given a per-query cost and a target \(E\rho^2\) (or \(\Phi\)), choose \((n_L, n_M, n_P, n_R)\) minimizing total queries \(n_L n_M n_P n_R\) subject to the reliability constraint. Because \(\sigma^2_\delta\) is monotone decreasing in each \(n\), the frontier is found by a direct grid search.

\subsection{Software and reproducibility}\label{software-and-reproducibility}

Models are fit with statsmodels 0.14 \texttt{MixedLM} using the crossed variance-components parameterization (a single constant grouping variable with one \texttt{vc\_formula} entry per facet). M1 is fit with the \texttt{bfgs} optimizer and M2 and M3 with \texttt{powell}, all at \texttt{maxiter\ =\ 3000}; every fit converged. The exact fitting script (\texttt{analysis-gtheory.py}) and its output (\texttt{results-gtheory.json}) are in the paper's repository folder, and the equivalent \texttt{lme4} and \texttt{statsmodels} specifications are given verbatim in Appendix A so the decomposition can be reproduced on any tidy iteration-level table with columns \texttt{brand,\ language,\ model,\ prompt\_id,\ iteration,\ outcome}. The response-level table for this corpus carries exactly those columns plus \texttt{is\_stability} and \texttt{prompt\_category} (Section 8).

\section{Results}\label{results}

This section reports the fitted decomposition. Every component below comes from a converged REML fit on the corpus of Section 3.1, reproducible from \texttt{analysis-gtheory.py} and its output \texttt{results-gtheory.json}. Point estimates are reported here; cluster-bootstrap 95\% confidence intervals on each component are deferred to the v2 finalization (Section 5.4, Limitation 1).

\subsection{Main-effects decomposition (full corpus)}\label{main-effects-decomposition-full-corpus}

Fitting M1 by REML on all 12,933 responses partitions the variance of a single response into four main-effect facets and a residual (Table 2). Query language is the largest systematic facet at 26.5\% of total variance; the residual bucket is larger at 69.3\% but conflates every interaction with pure within-prompt resampling and is decomposed on the stability subset (Section 5.2). Brand identity, the quantity a brand measurement exists to capture, is 1.5\%, an ICC of 0.000607 / 0.041502 = 0.0146. Model identity (1.6\%) and prompt (1.1\%) are of the same small order as brand. The residual is 69.3\%, but on the single-shot cells that bucket conflates every interaction with pure within-prompt resampling because M1 carries no cell term; M2 splits it. The reading of M1 is blunt: the variance of one AI answer about a brand is dominated by the residual bucket and by the language of the query, and brand identity contributes a fraction of one part in seventy. A single response carries almost no brand-discriminating signal.

\begin{table}[htbp]
\caption{M1, main-effects variance components, full corpus (n = 12,933). Outcome: response-level sentiment polarity.}
\small
\begin{tabularx}{\textwidth}{@{}>{\raggedright\arraybackslash}Xrr@{}}
\toprule
Source & Variance & ICC \\
\midrule

Query language & 0.010992 & 26.5\% \\
Model identity & 0.000681 & 1.6\% \\
Brand identity & 0.000607 & 1.5\% \\
Prompt & 0.000441 & 1.1\% \\
Residual (interactions + resampling) & 0.028781 & 69.3\% \\
\textbf{Total} & \textbf{0.041502} & 100\% \\
\bottomrule
\end{tabularx}
\end{table}

\emph{ICC of brand from a single response = 0.0146. Cluster-bootstrap 95\% CIs on every component are the v2 finalization (Section 5.4).}

\subsection{Isolating within-prompt resampling (Category D subset)}\label{isolating-within-prompt-resampling-category-d-subset}

M2 adds the full brand x language x model x prompt cell intercept and is fit on the 7,173 Category D responses, where cells are replicated. With the cell term present the residual becomes pure within-prompt resampling (Table 3). That residual is 34.8\% of variance: about a third of the movement in a replicated cell is the decoder sampling a different continuation on an identical prompt at temperature 0.3. The cell term itself is 29.6\%, and it is the brand-in-context interaction that M1 had lumped into its residual: the structure that says a specific brand, in a specific language, on a specific model and prompt, sits at a specific level. Query language rises to 32.0\% on this subset. Once the residual (pure within-prompt resampling, 34.8\%) and the brand-in-context cell interaction (29.6\%) are set aside, query language is the largest systematic between-condition facet, and the brand, model and prompt main effects stay small (1.6\%, 1.7\%, 0.3\%). Roughly, M1's 69.3\% residual splits into genuine cell-level interaction (about 30 points) and pure resampling (about 35 points); this split of the residual is identified on the Category D subset, and transporting it corpus-wide is an assumption.

\begin{table}[htbp]
\caption{M2, resampling-isolating variance components, Category D subset (n = 7,173). The residual is pure within-prompt resampling.}
\small
\begin{tabularx}{\textwidth}{@{}>{\raggedright\arraybackslash}Xrr@{}}
\toprule
Source & Variance & ICC \\
\midrule

Within-prompt resampling (residual) & 0.014873 & 34.8\% \\
Query language & 0.013681 & 32.0\% \\
Interactions (brand x language x model x prompt cell) & 0.012640 & 29.6\% \\
Model identity & 0.000733 & 1.7\% \\
Brand identity & 0.000698 & 1.6\% \\
Prompt (main effect) & 0.000118 & 0.3\% \\
\textbf{Total} & \textbf{0.042743} & 100\% \\
\bottomrule
\end{tabularx}
\end{table}

\emph{Validation:} the fitted resampling variance is 0.014873 (M2 residual). The direct model-free mean within-cell variance cross-check is computed and deposited with the bootstrap intervals in the v2 finalization (Section 5.4).

\subsection{The object-by-facet decomposition and the allocation frontier}\label{the-object-by-facet-decomposition-and-the-allocation-frontier}

M3 splits the M2 cell term into the object-by-facet interactions that matter for ranking brands (Table 4). Two results stand out. The brand's context-free true score, the pure brand main effect once interactions are removed, is 0.7\%, the smallest structural component in the model. And the only sizeable object-by-facet interaction is brand x language at 8.6\%: a brand's standing genuinely shifts with the language of the question, beyond what the language main effect already explains. This is the bilingual penalty expressed as a variance component. Brand x model and brand x prompt are both effectively zero, so a brand's relative standing is roughly stable across the three models and across paraphrases, but not across languages. The language main effect remains the largest systematic between-condition facet at 31.6\%, with pure within-prompt resampling (the residual) larger in absolute share at 34.8\%.

\begin{table}[htbp]
\caption{M3, object-by-facet variance components, Category D subset (n = 7,173).}
\small
\begin{tabularx}{\textwidth}{@{}>{\raggedright\arraybackslash}Xrr@{}}
\toprule
Source & Variance & Share \\
\midrule

Within-prompt resampling (residual) & 0.014873 & 34.8\% \\
Query language (main) & 0.013518 & 31.6\% \\
Higher-order interactions (cell) & 0.009522 & 22.3\% \\
Brand x language & 0.003664 & 8.6\% \\
Model identity (main) & 0.000736 & 1.7\% \\
Brand (object) & 0.000280 & 0.7\% \\
Prompt (main) & 0.000124 & 0.3\% \\
Brand x model & 0.000000 & 0.0\% \\
Brand x prompt & 0.000000 & 0.0\% \\
\textbf{Total} & \textbf{0.042717} & 100\% \\
\bottomrule
\end{tabularx}
\end{table}

Feeding the M3 components into the decision-study formula of Section 4.2 gives the marginal value of each way of spending the next block of queries, starting from the field-standard base design of one language, one model, five paraphrases, five repeats (25 queries, relative-error variance 0.006163). Table 5 reports it. Ranked by the variance reduction each block buys, languages come first (0.00462), then models (0.00167), then paraphrases (0.00125), and repeats last (0.00030): language diversity reduces relative-error variance about fifteen times as much as five more repeats. The ordering by reduction per query keeps language first (6.16 x 10\^{}-5) and repeats last (1.19 x 10\^{}-5), and there paraphrases edge models (5.00 x 10\^{}-5 against 3.33 x 10\^{}-5) because a model block adds two levels and costs twice the queries of a paraphrase block. Either way, the message is one message: repeats beyond five are the least efficient place to put a query. The field-standard habit of buying more repeats of the same prompt in the same language is the worst allocation of the budget.

\begin{table}[htbp]
\caption{Marginal reduction in relative-error variance from a single-language, single-model, five-paraphrase, five-repeat base design (1L x 1M x 5P x 5R = 25 queries, $\sigma^2_\delta$ = 0.006163). Computed from the M3 components.}
\footnotesize
\begin{tabularx}{\textwidth}{@{}>{\raggedright\arraybackslash}X>{\raggedright\arraybackslash}X>{\raggedleft\arraybackslash}Xrr>{\raggedleft\arraybackslash}X@{}}
\toprule
Where the queries go & New design & Extra queries & \(\sigma^2_\delta\) after & Reduction & Reduction per extra query \\
\midrule

5 more repeats & 1L x 1M x 5P x 10R & 25 & 0.005866 & 0.00030 & 1.19 x 10\^{}-5 \\
5 more paraphrases & 1L x 1M x 10P x 5R & 25 & 0.004914 & 0.00125 & 5.00 x 10\^{}-5 \\
2 more models & 1L x 3M x 5P x 5R & 50 & 0.004497 & 0.00167 & 3.33 x 10\^{}-5 \\
3 more languages & 4L x 1M x 5P x 5R & 75 & 0.001541 & 0.00462 & 6.16 x 10\^{}-5 \\
\bottomrule
\end{tabularx}
\end{table}

The precision-versus-budget frontier (Table 6) turns the same components into reliability. Two facts govern it. First, brand-ranking reliability is low everywhere. A single AI answer (one language, one model, one paraphrase, one repeat) has \(E\rho^2 \approx 0.01\), effectively no ability to rank brands from one screenshot. Hammering repeats at a single language and prompt reaches only 0.020 at twenty repeats, a twenty-fold budget for a gain of 0.01. Reliability climbs mainly by breadth: eight languages, one model, one paraphrase and ten repeats already reach \(E\rho^2 = 0.132\). Second, the design has a low ceiling. The full crossed 8-language, 3-model, 15-paraphrase design plateaus at \(E\rho^2 \approx 0.36\), and running it at 1, 5, 10 or 20 repeats moves the coefficient only from 0.347 to 0.365. Absolute-score reliability \(\Phi\) stays near 0.10 throughout, a direct consequence of the near-degenerate sentiment outcome, whose brand-object variance is 0.7\%. The non-monotone pair in the table is instructive: the 360-query design that buys paraphrase breadth (15 paraphrases, 1 repeat) reaches \(E\rho^2 = 0.347\), higher than the 600-query design that buys repeat depth (5 paraphrases, 5 repeats) at 0.332. Paraphrase breadth beats repeat depth even when it costs fewer queries.

\begin{table}[htbp]
\caption{Precision-versus-budget frontier. $E\rho^2$ is brand-ranking reliability, $\Phi$ is absolute-score reliability. Computed from the M3 components; cluster-bootstrap intervals are the v2 finalization.}
\footnotesize
\begin{tabular*}{\textwidth}{@{\extracolsep{\fill}}rrrrrrr}
\toprule
Languages & Models & Paraphrases & Repeats & Queries & \(E\rho^2\) & \(\Phi\) \\
\midrule

1 & 1 & 1 & 1 & 1 & 0.010 & 0.007 \\
1 & 1 & 1 & 20 & 20 & 0.020 & 0.010 \\
1 & 1 & 5 & 5 & 25 & 0.043 & 0.013 \\
1 & 3 & 5 & 5 & 75 & 0.059 & 0.015 \\
8 & 1 & 1 & 10 & 80 & 0.132 & 0.060 \\
8 & 3 & 15 & 1 & 360 & 0.347 & 0.102 \\
8 & 3 & 5 & 5 & 600 & 0.332 & 0.100 \\
8 & 3 & 15 & 5 & 1800 & 0.362 & 0.103 \\
8 & 3 & 15 & 10 & 3600 & 0.364 & 0.103 \\
8 & 3 & 15 & 20 & 7200 & 0.365 & 0.103 \\
\bottomrule
\end{tabular*}
\end{table}

\textbf{The allocation rule, stated as a procedure.} For a target \(E\rho^2 = t\), choose the least-cost \((n_L,n_M,n_P,n_R)\) on the grid with \(E\rho^2 \geq t\). On the fitted components the rule spends in the order languages, then models, then paraphrases, and reaches its ceiling on repeats first: past five repeats the constraint is never the binding one, so the optimizer never buys them. For any target reachable within the design (up to \(E\rho^2 \approx 0.36\)), the least-cost design maximizes language and model coverage before adding a sixth repeat. The structural expectation of Section 4.2, that repeats reach their reliability ceiling first, holds with a wide margin for this outcome.

\subsection{Further fits and confidence intervals still to run}\label{further-fits-and-confidence-intervals-still-to-run}

The primary decomposition (M1, M2, M3, the frontier, and the allocation) is computed and reported above. Four items remain and are named here so that the standing of each is explicit.

\begin{itemize}
\item
  \textbf{Cluster-bootstrap confidence intervals.} Resample the 20 brands with replacement, refit M1 to M3 on each replicate, and take percentile intervals on every variance component, every ICC, and the frontier coefficients. The direct model-free within-cell variance cross-check for the M2 residual (Section 5.2) is deposited with these intervals. Twenty brands is few, so the intervals will be wide; that width is the honest cost of generalizing to unseen brands. This is the v2 finalization.
\item
  \textbf{Recommendation-indicator outcome.} Refit M1, M2, M3 on a binary or ordinal recommendation outcome (brand named in the response, and its rank) derived from the stored response text via multilingual entity matching, using a crossed generalized linear mixed model (logit link) on the pooled iteration-level table; report the same components, ICCs, frontier and allocation. The brand-object signal is expected to be larger on this outcome than on the near-degenerate sentiment outcome, which would raise \(E\rho^2\) and \(\Phi\); the ordering of facet marginal values is a structural property that is not expected to change.
\item
  \textbf{Corporate-taxonomy corpus (Section 3.2).} Fit M1 to M3 on the 195-iteration, 6-brand corporate corpus and report whether the repeats-cap conclusion holds in a single-language, high-replication design.
\item
  \textbf{V-ZUG single-brand corpus (Section 3.3).} Fit the model and paraphrase facets on the 375-query V-ZUG corpus (single object, so brand becomes a fixed target and object variance is replaced by a category-visibility contrast) and locate it on the Table 6 frontier.
\end{itemize}

\section{Discussion}\label{discussion}

The decomposition converts a folk convention into a measured allocation. Three findings follow from the fits.

First, a single AI answer is almost useless for ranking brands. Brand identity accounts for 1.5\% of the variance of one response (ICC 0.0146) in the full-corpus fit, and the brand's context-free true score falls to 0.7\% once the object-by-facet model removes the interaction structure. The generalizability coefficient for ranking brands from a single one-language, one-model, one-paraphrase, one-repeat query is about 0.01. Whatever a single screenshot of an AI answer appears to say about where a brand stands, the measurement carries almost no brand-discriminating signal.

Second, the noise the field spends its budget on is the noise averaging already handles. Pure within-prompt resampling is 34.8\% of variance on the stability subset, but in an averaged design that term is divided by the entire query count, so it is the cheapest source to suppress. Adding a sixth through tenth repeat to a five-repeat design reduces relative-error variance by 0.0003, the smallest return of any facet. The structural prediction of Section 4.2, that the repeat facet has the fastest-decaying marginal value, holds in the fitted numbers, and the gap is wide: a language block buys about fifteen times the variance reduction of a repeat block from the same base design.

Third, the facet the field mostly ignores, query language, is the largest systematic between-condition source (larger than any other facet; only the residual resampling term, which is measurement noise rather than a between-condition source, has a bigger share on the stability subset). Language is 26.5\% of single-response variance in the full corpus and 31.6 to 32.0\% on the stability subset, and it carries a specific brand-level term: brand x language is 8.6\% of variance, a measurable bilingual penalty in which a brand's standing shifts with the language of the question beyond the language main effect. Brand x model and brand x prompt are near zero, so a brand's relative standing is roughly stable across models and paraphrases but not across languages. A measurement that never varies the query language is blind to the one interaction that most reorders a brand's score.

The structural argument stands behind all three. When a brand score is formed by averaging over a crossed design, the resampling term is divided by the entire query count while the language term is divided by the language count alone. For any positive components the marginal reliability gain from a repeat decays faster than the gain from a language, model, or paraphrase; the fitted magnitudes show the gap is decisive for this outcome. The precision frontier makes the ceiling explicit: even the full crossed 8-language, 3-model, 15-paraphrase design reaches only \(E\rho^2 \approx 0.36\) for ranking and \(\Phi \approx 0.10\) for absolute scoring, and running that design at 1, 5, 10 or 20 repeats moves \(E\rho^2\) only from 0.347 to 0.365. Reliability is bought by breadth across languages and models, not by depth of repetition. The low absolute ceiling is itself a finding about the outcome: sentiment polarity is 91.9\% exactly neutral, so the brand-object variance is small and \(\Phi\) cannot rise far. The recommendation-indicator refit (Section 5.4) is the test of whether a less degenerate outcome lifts the ceiling while preserving the facet ordering.

The framing clarifies what the two neighboring studies deliver and do not. A Gini coefficient reports that a category is concentrated; it does not report that the concentration estimate is itself unstable because it was measured in one language. A single-facet power tier reports that five iterations detect only large effects; it does not report where the next query is best spent when four facets compete. The variance-components view answers both, using standard generalizability theory, and it says the same thing twice: spend on languages and models, and stop buying repeats past five.

\section{Limitations}\label{limitations}

\begin{enumerate}
\def\labelenumi{\arabic{enumi}.}
\item
  \textbf{Results computed; confidence intervals pending.} The variance components, ICCs, precision frontier, and allocation in Section 5 are point estimates from converged REML fits. Cluster-bootstrap 95\% confidence intervals on each component, resampling the 20 brands, are not yet computed and are the v2 finalization; the point estimates should be read without interval guarantees until then, and the direct model-free within-cell cross-check of the resampling component (Section 5.2) is deposited with those intervals.
\item
  \textbf{Outcome.} The primary outcome is response-level sentiment polarity, which is strongly zero-inflated (median 0, mean -0.0205, 91.9\% of responses exactly neutral) and noisy in absolute terms on long list answers. A Gaussian mixed model on a near-degenerate outcome is a first-order approximation, and the brand-object signal is small on this outcome (0.7\% of variance in M3), which holds the absolute coefficient \(\Phi\) near 0.10 regardless of query budget. The recommendation-indicator refit (Section 5.4) is the needed cross-check; until it is run, the fitted ICCs are outcome-conditional.
\item
  \textbf{Low decoding temperature.} The three models were queried at temperature 0.3, which deliberately compresses within-prompt sampling variation. The resampling component estimated here (34.8\% of the stability-subset variance) is therefore a lower bound on what a default-temperature study would see, and the marginal value of repeats measured on this corpus does not transfer to studies run at higher temperature without adjustment.
\item
  \textbf{Mixed grounding across models.} Two of the three models (GPT-5.2, Gemini 3 Flash) were queried in parametric mode and one (Perplexity) in grounded retrieval mode. The model facet therefore confounds model identity with retrieval mode, and the model-by-brand and model-by-language interactions absorb any grounding effect. A design that crosses grounding with model would separate them; this corpus does not.
\item
  \textbf{Identifiability of low-order interactions.} With the full brand x language x model x prompt cell term in M3, the brand x prompt interaction is not cleanly separated from the cell term, so the allocation for the paraphrase facet may rest partly on the lumped cell term. This is visible in the fit: brand x prompt estimates to essentially zero while the cell term carries 22.3\%. A saturated model with every two- and three-way interaction estimated separately is the fuller specification; it is harder to fit stably and is deferred.
\item
  \textbf{Few facet levels.} Model has three levels and language eight, so their variance components rest on few degrees of freedom and their bootstrap intervals will be wide. The model component in particular should not be over-interpreted.
\item
  \textbf{Single corpus, single window.} All data come from one CEE corpus collected in one window in spring 2026 with three specific model versions. Model versions change and the components can drift with them. The two secondary corpora are described but not yet fit. External validity across markets, outcomes, and model generations is the open question, addressed by the further fits in Section 5.4.
\item
  \textbf{Cluster bootstrap granularity.} Confidence intervals resample 20 brands. Twenty is few, so the intervals will be wide, which is the honest consequence of generalizing to unseen brands from twenty.
\end{enumerate}

\section{Ethics and data availability}\label{ethics-and-data-availability}

\textbf{Ethics.} The study uses no human subjects and no personal data. It queries commercial LLM APIs about publicly traded or publicly known brands and analyzes the models' text. The only sensitive dimension is reputational: the analysis can surface that a model systematically scores a given brand lower in a given language. Such patterns are reported at the level of measurement error and cross-language divergence, not as claims about the brands' merits, and the release is of aggregate components and derived measures, not model outputs that could be quoted as authoritative statements about a company. Several vendors named in Section 2 compete with the author's company; their figures are attributed to them with sample and date and used only as independent evidence of the phenomenon this paper decomposes.

\textbf{Data availability.} Aggregate CEE data from the primary corpus (per-brand-per-language cross-language similarities, sentiment divergence, and stability scores) are part of the cross-language reputation dataset released under CC BY 4.0 at Zenodo, DOI 10.5281/zenodo.20794390. The response-level iteration table that is the direct input to the decomposition in this paper (columns \texttt{brand,\ language,\ model,\ prompt\_id,\ prompt\_category,\ iteration,\ is\_stability,\ sentiment\_score}) is available from the corresponding author on reasonable request. The analysis code that produces every component, ICC, the frontier, and the allocation is provided as \texttt{analysis-gtheory.py} with its output \texttt{results-gtheory.json} in the paper's repository folder, reproduced in Appendix A, and will be deposited with the Zenodo record. The category-ownership dataset used for the Section 2 comparison is at DOI 10.5281/zenodo.20788142, and the cross-market source-composition dataset referenced in the related work is at DOI 10.5281/zenodo.20829524; both are CC BY 4.0. Other open Rankfor datasets, the PersonaGen persona family, are hosted on the Hugging Face Hub under the Rankfor organization; they are separate resources and are not part of this corpus.

\textbf{Author contribution and funding.} Sole author. The work was conducted at Rankfor.AI with no external funding. The author declares a commercial interest: Rankfor.AI sells AI brand-visibility measurement, and the allocation rule in this paper is intended for internal use in sizing its measurements.

\begin{center}\rule{0.5\linewidth}{0.5pt}\end{center}

\section*{Appendix A. Runnable analysis}\label{appendix-a.-runnable-analysis}

The decomposition runs on any tidy iteration-level table with columns \texttt{brand,\ language,\ model,\ prompt\_id,\ iteration,\ y}, where \texttt{y} is the response-level outcome. For this corpus, \texttt{y} is \texttt{sentiment\_score} and the Category D subset is selected by \texttt{is\_stability}. The deposited script \texttt{analysis-gtheory.py} is the exact code that produced Tables 2 to 6 and \texttt{results-gtheory.json}; M1 is fit with \texttt{bfgs}, M2 and M3 with \texttt{powell}, all at \texttt{maxiter\ =\ 3000}.

\textbf{R (lme4).}

\begin{lstlisting}
library(lme4)
cee$cell <- with(cee, interaction(brand, language, model, prompt_id, drop = TRUE))
stab <- subset(cee, is_stability)

# M1  main-effects decomposition, full corpus
m1 <- lmer(y ~ 1 + (1|brand) + (1|language) + (1|model) + (1|prompt_id),
           data = cee, REML = TRUE)

# M2  resampling-isolating, Category D subset (residual = within-prompt resampling)
m2 <- lmer(y ~ 1 + (1|brand) + (1|language) + (1|model) + (1|prompt_id) + (1|cell),
           data = stab, REML = TRUE)

# M3  object-by-facet, Category D subset
m3 <- lmer(y ~ 1 + (1|brand) + (1|language) + (1|model) + (1|prompt_id) +
               (1|brand:language) + (1|brand:model) + (1|brand:prompt_id) + (1|cell),
           data = stab, REML = TRUE)
print(VarCorr(m3), comp = "Variance")
\end{lstlisting}

\textbf{Python (statsmodels).}

\begin{lstlisting}
import statsmodels.formula.api as smf
cee["grp"] = 1
vc_m3 = {
    "brand":       "0 + C(brand)",
    "language":    "0 + C(language)",
    "model":       "0 + C(model)",
    "prompt":      "0 + C(prompt_id)",
    "brandXlang":  "0 + C(brand):C(language)",
    "brandXmodel": "0 + C(brand):C(model)",
    "brandXprompt":"0 + C(brand):C(prompt_id)",
    "cell":        "0 + C(brand):C(language):C(model):C(prompt_id)",
}
stab = cee[cee.is_stability]
r = smf.mixedlm("y ~ 1", stab, groups=stab["grp"], vc_formula=vc_m3).fit(reml=True, method="powell", maxiter=3000)
comps = dict(zip(r.model.exog_vc.names, [float(x) for x in r.vcomp])); comps["resid"] = r.scale
total = sum(comps.values()); icc = {k: v/total for k, v in comps.items()}
\end{lstlisting}

\textbf{Decision study and allocation (from M3 components).}

\begin{lstlisting}
def rel_err(nL,nM,nP,nR,c):   # relative-error variance, ranking brands
    return (c["brandXlang"]/nL + c["brandXmodel"]/nM + c["brandXprompt"]/nP
            + c["cell"]/(nL*nM*nP) + c["resid"]/(nL*nM*nP*nR))
def Erho2(nL,nM,nP,nR,c):
    d = rel_err(nL,nM,nP,nR,c); return c["brand"]/(c["brand"]+d)
# allocation: least-cost design meeting a target reliability
best = min(((nL,nM,nP,nR) for nL in Ls for nM in Ms for nP in Ps for nR in Rs
            if Erho2(nL,nM,nP,nR,comps) >= target),
           key=lambda d: d[0]*d[1]*d[2]*d[3])
\end{lstlisting}


\begin{thebibliography}{99}
\bibitem{ref1} Atil, B., et al.~(2025). Non-Determinism of ``Deterministic'' LLM System Settings in Hosted Environments. In \emph{Proceedings of the 5th Workshop on Evaluation and Comparison of NLP Systems (Eval4NLP)}, EMNLP 2025. ACL Anthology 2025.eval4nlp-1.12 (arXiv:2408.04667).

\bibitem{ref2} Barbieri, F., Espinosa Anke, L., and Camacho-Collados, J. (2022). XLM-T: Multilingual language models in Twitter for sentiment analysis and beyond. In \emph{Proceedings of LREC 2022}, 258 to 266.

\bibitem{ref3} Bates, D., Mächler, M., Bolker, B., and Walker, S. (2015). Fitting linear mixed-effects models using lme4. \emph{Journal of Statistical Software}, 67(1), 1 to 48.

\bibitem{ref4} Brennan, R. L. (2001). \emph{Generalizability Theory}. Springer.

\bibitem{ref5} Conneau, A., Khandelwal, K., Goyal, N., Chaudhary, V., Wenzek, G., Guzmán, F., Grave, E., Ott, M., Zettlemoyer, L., and Stoyanov, V. (2020). Unsupervised cross-lingual representation learning at scale. In \emph{Proceedings of ACL 2020}, 8440 to 8451.

\bibitem{ref6} Cronbach, L. J., Gleser, G. C., Nanda, H., and Rajaratnam, N. (1972). \emph{The Dependability of Behavioral Measurements: Theory of Generalizability for Scores and Profiles}. Wiley.

\bibitem{ref7} Efron, B., and Tibshirani, R. J. (1993). \emph{An Introduction to the Bootstrap}. Chapman and Hall.

\bibitem{ref8} Graphite (2026). How many prompts and responses you need for a stable AI position estimate. Industry report, May 2026 (200 prompts by 400 responses). Referenced via Search Engine Land.

\bibitem{ref9} Indig, K. (2026). Measuring AI search: run-to-run citation stability. Referenced via Search Engine Land, June 2026.

\bibitem{ref10} Nakagawa, S., and Schielzeth, H. (2013). A general and simple method for obtaining R2 from generalized linear mixed-effects models. \emph{Methods in Ecology and Evolution}, 4(2), 133 to 142.

\bibitem{ref11} Ouyang, S., Zhang, J. M., Harman, M., and Wang, M. (2025). An empirical study of the non-determinism of ChatGPT in code generation. \emph{ACM Transactions on Software Engineering and Methodology} (arXiv:2308.02828).

\bibitem{ref12} Peec AI (2026). Fan-out and citation measurement reports. Industry reports, February to May 2026.

\bibitem{ref13} Profound (2025). Citation concentration in AI answers. Industry reports, June and November 2025.

\bibitem{ref14} Searle, S. R., Casella, G., and McCulloch, C. E. (1992). \emph{Variance Components}. Wiley.

\bibitem{ref15} Seabold, S., and Perktold, J. (2010). statsmodels: Econometric and statistical modeling with Python. In \emph{Proceedings of the 9th Python in Science Conference}.

\bibitem{ref16} Shavelson, R. J., and Webb, N. M. (1991). \emph{Generalizability Theory: A Primer}. Sage.

\bibitem{ref17} Understanding and mitigating numerical sources of nondeterminism in LLM inference (2025). arXiv:2506.09501.

\bibitem{ref18} Zatuchin, D. (2026a). The category-ownership map: measuring which brands own which recommendation categories in AI. Dataset DOI 10.5281/zenodo.20788142. Preprint; under review at the Journal of Marketing Analytics.

\bibitem{ref19} Zatuchin, D. (2026b). The dice-roll method: power and generalizability for measuring LLM brand recommendations. Preprint; under peer review.

\bibitem{ref20} Zatuchin, D. (2026c). Cross-language AI brand reputation across European languages, including the Central and Eastern European arm used here. Dataset DOI 10.5281/zenodo.20794390. Preprint.

\bibitem{ref21} Zatuchin, D. (2026d). Six corporates, three visibility patterns: a dice-roll audit of large-cap AI visibility. Rankfor.AI research article, open.rankfor.ai, 2026.

\bibitem{ref22} Zatuchin, D. (2026e). V-ZUG and the incumbent shadow: unbranded-query visibility in premium appliances. Rankfor.AI research article, open.rankfor.ai, 2026.
\end{thebibliography}

\end{document}
